% (C) 2026 Ralf Stephan, in collaboration with Claude Code.
% Released under CC0 1.0 Universal (public-domain dedication).
% See https://creativecommons.org/publicdomain/zero/1.0/

\documentclass[11pt]{article}
\usepackage[T1]{fontenc}
\usepackage[utf8]{inputenc}
\usepackage{lmodern}
\usepackage[margin=2.5cm]{geometry}
\usepackage{amsmath,amssymb}
\usepackage{braket}
\usepackage[hidelinks]{hyperref}

% Lean names and paths: verbatim, breakable at . _ /
\DeclareUrlCommand\lean{\urlstyle{tt}}
\def\UrlBreaks{\do\.\do\_\do\/}
% Extra stretch for paragraphs with long, unbreakable Lean names
\emergencystretch=6em

\title{Open problems:\\ Inconsistent convergence in Hilbert space}
\author{}
\date{}

\begin{document}
\maketitle

\noindent
Open problems raised by the technical brief in this directory, and by its Lean formalization
(see \lean{formalization-log.tex}). Items marked \emph{research} are open research questions,
not needed to check the brief.

\section*{From the brief}
\begin{enumerate}
\item \textbf{A criterion of physical convergence.} \emph{Research.} The brief concludes that the $\ell^2$ topology
  ``is too weak to take into account all conditions for a meaningful criterion of physical
  convergence''. Which topology is the right one, and whether it is the Schwartz topology, is
  left open. (Brief, Sec.~4.)
\end{enumerate}

\end{document}
